% Generated draft: requires mathematical and source review.
\documentclass[11pt]{article}
\usepackage[margin=25mm]{geometry}
\usepackage{fontspec}
\setmainfont{Latin Modern Roman}
\usepackage{amsmath,amssymb,mathtools,unicode-math}
\setmathfont{Latin Modern Math}
\usepackage{enumitem,needspace,xurl,hyperref}
\setlength{\parindent}{0pt}
\setlength{\parskip}{5pt}
\setlength{\emergencystretch}{3em}
\newcommand{\N}{\mathbb{N}}
\newcommand{\Z}{\mathbb{Z}}
\newcommand{\Q}{\mathbb{Q}}
\newcommand{\R}{\mathbb{R}}
\newcommand{\C}{\mathbb{C}}
\newcommand{\F}{\mathbb{F}}
\newcommand{\A}{\mathbb{A}}
\newcommand{\PP}{\mathbb{P}}
\newcommand{\E}{\mathbb{E}}
\newcommand{\eps}{\varepsilon}
\newcommand{\dd}{\,\mathrm{d}}
\begin{document}
\section*{UnsolvedMath ID: 30006503}
\textbf{Source problem code:} catalogue entry
\subsection{Definitions and mathematical statement}
A Banach space \(X\) has the unconditional martingale differences property (UMD) if, for some \(q\in(1,\infty)\), there is a constant \(\beta\) such that every finite \(X\)-valued martingale difference sequence \((d_k)_{k=1}^n\) in \(L^q\), on any probability space, and every choice \(\varepsilon_k\in\{-1,1\}\) satisfy
\[
\left\|\sum_{k=1}^n\varepsilon_k d_k\right\|_{L^q(X)}
\leq \beta\left\|\sum_{k=1}^n d_k\right\|_{L^q(X)}.
\]
For \(0<\theta<\pi\), write
\[
\Sigma_\theta=\{z\in\C\setminus\{0\}:|\arg z|<\theta\},\qquad
H^\infty(\Sigma_\theta)=\{f:\Sigma_\theta\to\C:f\text{ is holomorphic and bounded}\},
\]
with norm \(\|f\|_\infty=\sup_{z\in\Sigma_\theta}|f(z)|\). In the usual sectorial-calculus formulation, bounded \(H^\infty(\Sigma_\theta)\) calculus requires sectoriality of angle strictly less than \(\theta\) and an estimate
\[
\|f(A)\|_{\mathcal L(E)}\leq C\|f\|_\infty
\qquad(f\in H_0^\infty(\Sigma_\theta)),
\]
where \(E\) is the underlying complex Banach space and \(H_0^\infty(\Sigma_\theta)\) consists of functions in \(H^\infty(\Sigma_\theta)\) for which there exist \(c,\alpha>0\) such that
\[
|f(z)|\leq c\min\{|z|^\alpha,|z|^{-\alpha}\}\qquad(z\in\Sigma_\theta).
\]
Here \(f(A)\) is defined by the sectorial holomorphic functional calculus.

\textbf{Provisional statement, retaining the supplied quantifiers.} For every UMD Banach space \(X\) and every \(p\in(1,\infty)\), does the Ornstein--Uhlenbeck operator \(A=-L_X\), in the realization intended by the source, have a bounded \(H^\infty(\Sigma_\theta)\) functional calculus for every \(0<\theta<\theta_p\)? The bound may depend on \(X,p,\theta\); the cleaned statement imposes no explicit uniformity in these parameters.

\textbf{Unresolved qualifications.} The record does not define \(\theta_p\), the underlying Gaussian measure space, the ambient function space, the domain or normalization of \(L_X\), or whether a dimension is fixed or quantified. It also does not specify whether \(X\) is complex or how a real space is to be complexified. Thus the displayed question preserves the source's symbolic formulation but does not determine a fully specified operator problem. The sector definition above applies to angles in \((0,\pi)\); the missing definition of \(\theta_p\) must establish the intended angular range. The direction \(0<\theta<\theta_p\) is retained exactly and requires source review: it asks for calculus at arbitrarily small positive sector angles. The original extraction additionally appears to ask for a dimension-uniform bound over \(d\geq1\), involving an undefined quantity denoted by a damaged expression containing \(M,p,d,\theta_p\), and \(X\). Its notation and quantifiers cannot be recovered reliably from this record, so that additional target remains unresolved.
\subsection{Short English statement}
The supplied question asks whether every UMD Banach space has the stated Ornstein--Uhlenbeck functional calculus for every \(1<p<\infty\) and every \(0<\theta<\theta_p\). The threshold, operator realization, and an apparent stronger dimension-uniform question are unspecified. This is a provisional statement requiring human mathematical review.
\subsection{Sources}
\begin{itemize}
\item Supplied local UnsolvedMath registry record, permanent ID \(30006503\), title \emph{Optimal Sector Angles for UMD Ornstein--Uhlenbeck Calculus}; exact source problem code \texttt{catalogue entry}. The record lists proposed year 2025 and gives no individual proposer.
\item Supplied source citation: \emph{Recent Developments in SPDEs and BSDEs meet Harmonic and Functional Analysis (2026)}, Oberwolfach Reports, pp.\ 2843--2916, DOI \url{https://doi.org/10.4171/owr/2025/53}. The literature list labels this \emph{Oberwolfach Report 53/2025}. The supplied title year 2026 and report/proposed year 2025 are retained without reconciliation. No precise problem page or speaker attribution is supplied.
\item Supplied literature reference: \emph{Holomorphic functional calculus and vector-valued Littlewood--Paley theory}, \url{https://doi.org/10.4171/JEMS/1430}. Authors and a theorem locating the claimed implication are not supplied.
\item Literature assessment dated \texttt{2026-08-22}: the record labels the problem \texttt{partially\_solved} and attributes to Kalton--Weis theory bounded \(H^\infty\) calculus for some angle below what it calls the optimal scalar angle, while describing the assertion for all angles below the threshold as open. This is supplied dated curation, not an independently verified literature assessment. The background expressly requests independent verification of time-sensitive status and cited claims.
\item The registry labels the statement \texttt{corrected\_verified}, explaining that only formula delimiters, notation, grammar, and typography were repaired and that mathematical content and scope were preserved. This supplied editorial label does not establish mathematical correctness or resolve the missing definitions.
\end{itemize}
\end{document}
